\section{Proof of the Conway criterion}
\label{sec:mainproof}

We can now assemble the preceding lemmas.

\begin{proof}[Proof of \cref{thm:mainPL}]
The period--area identity \cref{thm:periodarea} gives
\[
 \det(t,u)=2\Area(T)>0,
\]
which proves part~(a) and makes \(L=\Z t+\Z u\) a rank-two lattice.

Project the two representatives \(T\) and \(R(T)\) to \(\Torus=\R^2/L\).  By \cref{lem:boundaryzero}, their sum is a boundaryless nonnegative two-chain.  By the finite common-subdivision and constant-multiplicity argument of \cref{sec:torus}, this chain has constant multiplicity \(k\).  Exact area comparison and positivity cancel the lattice determinant and give \(k=1\).  Therefore \cref{thm:toruspartition} supplies a multiplicity-one partition of the torus by the two orbit representatives.

Take the full inverse image under the quotient map \(\pi:\R^2\to\Torus\).  The inverse image of \(\pi(T)\) is exactly the collection \(T+mt+nu\), and the inverse image of \(\pi(R(T))\) is exactly \(R(T)+mt+nu\), for \((m,n)\in\Z^2\).  Since the torus is covered, the plane is covered.  Since torus interiors are disjoint, the interiors of distinct lifted tiles are disjoint.  This proves parts~(b) and~(c).

Finally, \cref{prop:edgepaired} identifies every neighboring pair by the translated edge equation or one of the four designated half-turn equations.  The periods themselves are compositions of those half-turns by \cref{eq:fourcompositions}, so every placement uses only translations and half-turns generated by the Conway data.  This proves part~(d) and the periodic-tiling conclusion.
\end{proof}

\subsection{Degenerate marked points and empty arcs}

The proof does not require the six marked points to be pairwise distinct.  If, say, \(B=C\), then \(b\) may be the singleton path \((B)\); its half-turn equation is trivial and \(\PathJ(b)=0=\det(B,B)\).  All affine and chain identities remain valid.  Likewise \(A=B\) makes \(a\) empty and forces \(D=E\) through the translated-pair condition.

The usual requirement that at least three marked points be distinct prevents vacuous labellings in the topological statement.  In the certified PL theorem, positive disk area already forces \(\det(t,u)>0\), hence enough geometric nondegeneracy for the construction.  Keeping the conventional requirement in \code{ConwayOK} makes the arithmetic statement match the established criterion.

\subsection{Why the proof is stronger than a drawing}

There are three logically separate tasks:
\begin{enumerate}[label=(\roman*)]
\item construct candidate placements;
\item show every open boundary arc has a matching neighbor; and
\item prove global coverage and nonoverlap, including at vertices.
\end{enumerate}
The strip argument directly addresses the first two.  The torus cycle and area computation address the third without a case analysis of the possible coincidences among \(A,\ldots,F\).  In particular, no angle-sum calculation is required.  All vertex configurations are absorbed into the single equality \(k=1\).

\subsection{A bounded-patch corollary}

For formal arithmetic it is useful to replace the phrase ``the infinite plane is tiled'' by a family of finite statements.

\begin{corollary}[Primitive-recursive bounded patches]
\label{cor:patches}
Given a normalized PL Conway datum and a natural radius \(N\), one can compute finite bounds \(M(N)\) such that all tiles meeting the square \([-N,N]^2\) occur among
\[
 T^i_{m,n},
 \qquad i\in\{0,1\},
 \qquad |m|,|n|\leq M(N).
\]
The resulting finite list covers the square and has pairwise disjoint interiors there.
\end{corollary}
\begin{proof}
Use \cref{eq:latticecoords}.  The lattice coordinates of every vertex of either orbit representative have finite rational minima and maxima.  If a translated tile meets the query square, the integer lattice indices must lie in intervals obtained by subtracting those extrema from the extrema of the square's lattice coordinates.  Taking floors and ceilings gives a primitive-recursive bound.  The coverage and disjointness assertions are restrictions of the global tiling just proved; alternatively they can be checked directly by clipping to the query square.
\end{proof}

This corollary is one convenient arithmetic reading of the global theorem.  An even smaller certificate is the multiplicity-one torus partition, because periodicity lets every bounded patch be reconstructed from it.
